% !PRINTED 612
\weberpage{612}

%<<CONT
je nachdem das Ideal $\mathfrak a$ in der Klasse $A_i$ oder in einer anderen
Klasse enthalten ist. Multiplizieren wir also die Formel (2) mit
$\chi_\varkappa(\overline{A_i'})$ und summieren in bezug auf $\varkappa$, so folgt:

\weq{3}{
\frac{1}{h}\overset{\varkappa}{\sum}\chi_\varkappa(A_i') \log Q_k
% [sic] printed $\log Q_k$ with Latin $k$, but the summation index is $\varkappa$, so $\log Q_\varkappa$ is meant
= \overset{(1)}{\sum}\frac{1}{N(\mathfrak p)^s} + \frac{1}{2}\overset{(2)}{\sum}\frac{1}{N(\mathfrak p)^{2s}} + \frac{1}{3}\overset{(3)}{\sum}\frac{1}{N(\mathfrak p)^{3s}} + \cdots,
}
worin sich die Summen $\overset{(1)}{\sum}, \overset{(2)}{\sum}, \overset{(3)}{\sum}, \ldots$ auf alle Primideale $\mathfrak p$ erstrecken,
deren erste, zweite, dritte, $\ldots$ Potenz in der Klasse $\overline{A_i}$ enthalten ist.

Wir gehen in (3) zur Grenze $s = 1$ über; die Summen $\overset{(2)}{\sum}, \overset{(3)}{\sum}, \ldots$
behalten dabei einen endlichen Wert, da ihre unbedingte Konvergenz für $s = 1$ nicht aufhört. Dasselbe gilt von dem Teil
der Summe $\overset{(1)}{\sum}$, der sich auf Primideale von höherem als dem
ersten Grade bezieht, und es bleibt also nur noch fraglich, ob
die Summe

\[
\sum \frac{1}{N(\mathfrak p)^s},
\]

erstreckt über alle Primideale ersten Grades der Klasse $\overline{A_i}$, endlich
bleibt oder unendlich wird. Ersteres würde der Fall sein, wenn
es in $\overline{A_i}$ gar keine oder nur eine endliche Anzahl von Primidealen
ersten Grades gäbe.

Setzen wir aber die Existenz eines Klassenkörpers vom Grade $h$ % faint pre-$h$ mark is scan/print damage, not a character
voraus, so folgt aus § 164, 8. und § 163, (21), daß

\[
\log Q_2,\ \log Q_3,\ \ldots,\ \log Q_h
\]

endliche Grenzwerte behalten, während $\log Q_1$ unendlich wird. Die
linke Seite von (3) wird also unendlich und die rechte Seite kann
nicht endlich bleiben.

Damit ist bewiesen:

12. \gesperrt{Wenn ein Klassenkörper existiert, dessen Grad
nicht höher ist als die Klassenzahl, so enthält jede
Klasse unendlich viele Primideale ersten Grades.}

\begin{center}
§ 166. \textbf{\gesperrt{Primideale in den Idealklassen.}}
\end{center}

Nehmen wir als Gruppe zunächst die Gesamtheit der Zahlen $O$
des Körpers $\Omega$ ohne die Null, so sind die Klassen die Idealklassen
 des Körpers $\Omega$, die, wie wir gesehen haben, den Klassen
quadratischer Formen der Diskriminante $\varDelta$ eindeutig zugeordnet
%CONT>>
